\section{Benchmark Design}
\label{sec:design}

We adopt the four design principles SeBS states --- relevance, usability, a
scientific evaluation methodology, and extensibility~\cite{sebs2021} --- and
add one the sandbox setting forces on us: \emph{capability honesty}. Providers
do not offer the same operations, and a suite that scores a missing operation
as zero produces a ranking that rewards having fewer features.

\subsection{Five dimensions, two scales}

Table~\ref{tab:dimensions} is the specification. The rows are dimensions: what
system capability is being measured. The columns are scales: one sandbox, or
many at once.

\begin{table}[t]
\caption{The benchmark specification. Scalability is not a sixth row: it is the
cluster column of the other five.}
\label{tab:dimensions}
\small
\begin{tabular}{@{}lll@{}}
\toprule
Dimension & Single sandbox & Cluster \\
\midrule
Lifecycle & create-to-first-exec & ready throughput, burst p99 \\
Execution & per-command latency & concurrent command throughput \\
Compute   & build completion time & concurrent build slowdown \\
Storage   & metadata ops/s, IOPS & concurrent I/O slowdown \\
Network   & egress RTT, throughput & concurrent network slowdown \\
\bottomrule
\end{tabular}
\end{table}

Two choices in that table need defending.

\paragraph{Scalability is not a dimension.} An earlier version of this design
listed scalability beside lifecycle, compute and storage. That was wrong.
Scalability is not an independent resource; it is what happens to the other
five under load. Keeping it as a peer row would score the same underlying
behaviour twice --- once as ``storage'' and again as ``scalability'' ---
and any composite ranking built on the result would double-count it.

\paragraph{Orthogonal does not mean statistically independent.} The five
dimensions have distinct \emph{measurement targets}, not uncorrelated outputs.
A build time depends on CPU, filesystem and memory together. What the design
requires is weaker and achievable: each dimension measures a different thing,
each microbenchmark isolates one dominant bottleneck, each result answers a
question no other result answers, and no single underlying behaviour is scored
more than once.

\subsection{Microbenchmarks}

Eight microbenchmarks populate the table. Each exists to explain a number in
the application results, not to be a ranking of its own.

\begin{description}
  \item[B01 Cold create.] Create a sandbox and run the first command.
  Measured to the completion of that command, not to the control plane
  reporting ``running'': a provider that returns early has moved its startup
  cost into the user's first command rather than removed it.
  \item[B02 Resume.] Write a sentinel, suspend, resume, run a command, read
  the sentinel back. Reports resume-to-ready latency \emph{and} a state
  correctness rate.
  \item[B03 Lifecycle churn.] Create and destroy continuously. Reports
  latency drift and cleanup failures; a provider that leaks sandboxes after a
  burst is charging for them.
  \item[B04 Short command.] A no-op command on a live sandbox, repeated. This
  measures the exec transport, which for an interactive session is the cost
  that is paid $n$ times.
  \item[B06 CPU build.] Compile and run a fixed C program written by the
  benchmark itself, so the measurement never contains a package download and
  never shifts when an upstream toolchain does.
  \item[B09 Small files.] Create, stat and delete 5{,}000 small files. An
  unpacked dependency tree is a metadata workload, not a bandwidth one, so
  this is the storage number that predicts \texttt{npm ci}.
  \item[B12 Egress.] Fetch a fixed endpoint from inside the sandbox, timed by
  \texttt{curl} so the exec transport is excluded. Both the endpoint and the
  client location are part of the result; change either and the number stops
  comparing.
  \item[B14--B16 Cluster.] Burst provisioning, sustained open-loop
  provisioning, and noisy-neighbour isolation (\S\ref{sec:design:cluster}).
\end{description}

Four candidate microbenchmarks were cut. Browser automation and data
processing were dropped because both re-measure compute, storage and network
without isolating anything; parallel execution and long-running execution were
dropped as benchmarks and kept as \emph{parameters}, since the same workload
run at higher concurrency or for longer is not a different benchmark. Model
inference and model training were cut because they would make the attached
accelerator, not the sandbox, the thing being ranked.

\subsection{Four application workloads}

Application workloads describe what a developer or an agent waits for. They
mix dimensions on purpose and are reported in their own table.

\begin{description}
  \item[A. Environment preparation.] \texttt{git clone} at a pinned commit,
  then a dependency install. Dominated by network and filesystem metadata.
  \item[B. Code build and execution.] Build, then test. Dominated by CPU,
  memory and disk.
  \item[C. Interactive development.] Start a dev server, wait for it to serve,
  change a file, wait for the reload. Exercises long-lived processes, file
  watching and port exposure.
  \item[D. Stateful continuation.] Write work, suspend, resume, verify. The
  only application workload with a correctness result as well as a latency.
\end{description}

No application workload calls a language model. Model latency and output
non-determinism would be added to a measurement of infrastructure, and would
dominate it.

\subsection{The reinforcement-learning profile}

The RL profile reuses the dimensions above rather than adding new ones: a
reset is lifecycle plus storage, a step is execution, a rollout is cluster
scale. It contributes three benchmarks.

\begin{description}
  \item[Environment step.] Action to observation on a live environment, with
  fixed step compute and a configurable observation size, so what the series
  measures is the sandbox's own scheduling and transport overhead.
  \item[Environment reset.] Run three times, once per mechanism ---
  in-process, sandbox recreate, snapshot restore --- and reported as three
  results. An environment reset is not necessarily a sandbox restart, and
  collapsing the three would compare incomparable operations.
  \item[Parallel rollouts.] $N$ environments concurrently, each running a
  fixed episode. Reports transitions per second and the fraction of episodes
  that completed intact, because a rollout that breaks partway costs the
  trainer the whole episode.
\end{description}

\subsection{Lifecycle profiles}

The same workload is driven under four lifecycle patterns, which is why
duration and concurrency are run parameters rather than separate benchmarks:
\emph{ephemeral} (create, run once, destroy), \emph{persistent} (create, run
many times, destroy), \emph{recoverable} (create, run, checkpoint, resume) and
\emph{pooled} (prewarm, acquire, run, release).

\subsection{Cluster experiments}
\label{sec:design:cluster}

Three cluster metrics carry the scale story.

\paragraph{Sustainable provisioning throughput.} Not creation requests
accepted per second, but sandboxes that reach their first command per second,
under a stated service-level objective:
\[
  T_{\mathrm{sustained}} =
  \frac{\text{sandboxes that reached first exec}}{\text{elapsed seconds}}
  \;\Big|\; p_{99} \le L,\; \text{success} \ge S.
\]
We report $L$ and $S$ with every value. A control plane that accepts a
thousand creates and queues them has not provisioned a thousand sandboxes.

\paragraph{Effective density.} Not how many sandboxes a host admits, but how
many run SLO-compliant work:
\[
  D_{\mathrm{effective}} =
  \frac{\text{SLO-compliant active sandboxes}}{\text{physical hosts}}.
\]
For hosted providers the denominator is unobservable, so this metric applies
only to the self-hosted infrastructure configuration
(\S\ref{sec:methodology}).

\paragraph{Performance isolation.} Fixed work measured alone and then beside
$N$ interferers:
\[
  \text{slowdown} = \frac{t_{\mathrm{contended}}}{t_{\mathrm{baseline}}} - 1 .
\]

\subsection{Result status}

Every benchmark against every provider ends in exactly one of four states, and
all four are printed in every table.

\begin{description}
  \item[completed] The benchmark ran; its numbers are reportable.
  \item[unsupported] The provider offers no such capability. This is a
  statement about features, not about speed.
  \item[failed] The provider supports the operation and the run broke. The
  reason is printed in full.
  \item[skipped] The configuration did not request it.
\end{description}

The distinction is not cosmetic. Collapsing \emph{unsupported} into a zero
score makes a provider look slow for a feature it never claimed, and makes the
resulting leaderboard reward a smaller feature set.
